% generated from Paper 9/anccft9.tex
\chapter{Algorithmic non-commutative class field theory, IX: exceptional zeros, the graph $\cL$-invariant, and the operator-algebraic
Mazur--Tate--Teitelbaum formula}\label{chap:IX}
\chaptermark{IX}
\begin{chapterabstract}
Section~3 of [Ch.~\ref{chap:III}] recorded, without proof, that the cleared voltage pencil
$g(T)=T^{2}h(T)$ of an abelian $\Z/\ell^{k}$ tower over a $(q{+}1)$-regular graph has a
double zero at the trivial character, and that the constant term of the Iwasawa growth
law obeys $\nu_0=\ord_\ell\kappa(Y_0)-\mu+\sum_j\varepsilon_j$ in all four towers of
[Ch.~\ref{chap:II}, Table~2]. This paper turns both observations into theorems and identifies the
number they conceal. The double zero is forced by the symmetry $D(t^{-1})=D(t)^{t}$ of
the pencil; its second-order coefficient has the closed form
\[
\tfrac12\Theta''(0)\;=\;-\,\kappa(Y_0)\,\|h_\alpha\|^{2},
\]
where $h_\alpha$ is the harmonic representative of the voltage cochain $\alpha$ --- its
orthogonal projection onto the cycle space, equivalently onto the kernel of the
Hecke--Dirac operator of [Ch.~\ref{chap:II}] --- and $\kappa(Y_0)=|\Jac(Y_0)|=\frac1n\prod_f\#E_f(\F_p)$
is the algebraic $L$-value of [Ch.~\ref{chap:I}]. We therefore define the graph $\cL$-invariant
$\cL(Y_\bullet):=\tfrac12\Theta''(0)/\kappa(Y_0)=-\|h_\alpha\|^{2}$, a rational number
depending only on the class $[\alpha]\in H^{1}(Y_0;\Z)$ and not on $\ell$; for a
single-chord voltage, $-\tfrac12\Theta''(0)$ counts the spanning trees avoiding the
chord. The bookkeeping identity of [Ch.~\ref{chap:III}] becomes a theorem with an exact description of
the finitely supported defects $\varepsilon_j$, and the specification's proposed
alignment is replaced by the correct one: in the exceptional-zero regime $\lambda_h=0$
all defects vanish and $\nu_0=-\ord_\ell\cL(Y_\bullet)$, while in general
$\nu_0+\ord_\ell\cL(Y_\bullet)=(\ord_\ell a_2-\mu)+\sum_j\varepsilon_j$. On the
operator-algebraic side we prove a twisted Bass identity for the character-twisted
Satake dynamics of [Ch.~\ref{chap:I}] and show that at the critical inverse temperature $\beta=1$ the
twisted Bass determinant equals $q^{-n}\Theta(t-1)$: its zero at the trivial character
is simple in the temperature direction with derivative proportional to $\kappa(Y_0)$,
of order two in the character direction with second derivative proportional to
$\cL(Y_\bullet)\kappa(Y_0)$, and the partition function of the critical KMS family has a
double pole with leading coefficient $n(q-q^{-1})/\cL(Y_\bullet)$ --- the
Mazur--Tate--Teitelbaum and Greenberg--Stevens shapes, with the roles of algebraic
$L$-value and $\cL$-invariant made explicit. The abelian shadow limit of [Ch.~\ref{chap:VIII}] carries
$\Theta$ as the characteristic element of its equivariant $K_0$. For the non-normal
Iwahori tower there is no character direction: its zero is simple at every level, and
the constancy of its residue $n_k\kappa(Y_k)q^{-n_k}$ along the tower is exactly the
pseudo-Iwasawa law, with the Euler term $-k$ appearing as the $p$-adic valuation of the
vertex count. No $\cL$-invariant exists for that tower, and we say so.
\end{chapterabstract}

\section*{Status of the results}

Every numbered statement is proved. The proofs use only results of Papers I--VIII,
finite-dimensional linear algebra (a Schur complement replaces analytic perturbation
theory), the Weierstrass preparation theorem for $\Zl[[T]]$ \cite{Wa}, and the character
product formula for abelian covers already used in [Ch.~\ref{chap:II}, \S5] and [Ch.~\ref{chap:IV}]. The twisted Bass
identity is the Artin--Ihara $L$-function formula of Stark--Terras \cite{ST}; we give a
six-line proof so that nothing is quoted. Every identity was machine-verified exactly
on the four abelian towers of [Ch.~\ref{chap:II}, Table~2] --- growth laws recomputed directly from the
reduced Laplacians of the derived graphs up to $128$ vertices --- and on the Iwahori tower
at levels $k\le3$ (Table~\ref{IX:tab:battery}). The Mazur--Tate--Teitelbaum \cite{MTT} and
Greenberg--Stevens \cite{GS} formulas are cited for the shape of the analogy only;
nothing here depends on them. Section~\ref{IX:sec:scope} records what is open and five
corrections to the specification that prompted this paper.

\section{The exceptional zero of the voltage pencil}\label{IX:sec:zero}

Setting as in [Ch.~\ref{chap:II}, \S5] and [Ch.~\ref{chap:IV}]. $Y_0$ is a finite connected $(q{+}1)$-regular graph
with $n$ vertices and $m=\frac{(q+1)n}2$ edges, $r=m-n+1$; fix an orientation $E^+$ of
the edges and a voltage $\alpha\colon E^+\to\Z$, extended to oriented edges by
$\alpha(\bar e)=-\alpha(e)$, so that $\alpha$ is an integer $1$-cochain. The voltage
adjacency and Laplacian pencils are
\[
A(t)_{uv}=\sum_{e\colon u\to v}t^{\alpha(e)},\qquad
D(t)=(q{+}1)I_n-A(t)\ \in M_n\bigl(\Z[t,t^{-1}]\bigr),
\]
and $\Theta(T):=\det D(1+T)\in\Z[T,(1+T)^{-1}]\subset\La:=\Zl[[T]]$ is the graph
$\ell$-adic $L$-function, the characteristic element of the limit module of
[Ch.~\ref{chap:IV}, Thm.~2.1]. Its cleared form $g(T)=(1+T)^{d}\Theta(T)\in\Z[T]$, $d$ minimal, is the
polynomial of [Ch.~\ref{chap:II}, Rem.~5.4] and [Ch.~\ref{chap:III}, \S3]; since $(1+T)^{d}=1+dT+\cdots$ and
$\Theta(0)=\Theta'(0)=0$ (Theorem~\ref{IX:thm:zero}), the coefficients
$a_i$ of $T^{i}$ in $\Theta$ and in $g$ agree for $i\le2$. We write
$\kappa_0:=\kappa(Y_0)=|\Jac(Y_0)|=\frac1n\prod_f\#E_f(\F_p)$, the algebraic $L$-value
of [Ch.~\ref{chap:I}, Thm.~3.5], and $\Delta_0=D(1)$ for the Laplacian of $Y_0$.

Hodge data. Let $\partial\colon\Z^{E^+}\to\Z^{V}$, $\partial e=t(e)-o(e)$, be the
boundary map, so $\Delta_0=\partial\partial^{t}$ [Ch.~\ref{chap:II}, Lem.~3.2]; the cycle space
$\ker\partial=H_1(Y_0;\Z)$ and the cut space $\operatorname{im}\partial^{t}$ are
orthogonal complements in $\Q^{E^+}$. With $\Delta_0^{+}$ the Moore--Penrose inverse
(the inverse of $\Delta_0$ on $\mathbf1^{\perp}$, zero on $\mathbf1$), the orthogonal
projection onto the cycle space is
\[
\Pi:=I-\partial^{t}\Delta_0^{+}\partial ,
\]
which is also the projection onto $\ker\mathsf D_p\cap\ell^{2}(E^+)$ for the
Hecke--Dirac operator $\mathsf D_p$ of [Ch.~\ref{chap:II}, Def.~3.1]. Put $h_\alpha:=\Pi\alpha$, the
harmonic representative of $\alpha$, and $\|h_\alpha\|^{2}=\alpha^{t}\Pi\alpha
=\|\alpha\|^{2}-(\partial\alpha)^{t}\Delta_0^{+}(\partial\alpha)\in\Q_{\ge0}$.

\begin{lemma}[Symmetry and switching]\label{IX:lem:sym}
$D(t^{-1})=D(t)^{t}$, hence $f(t):=\det D(t)$ satisfies $f(t^{-1})=f(t)$. For
$\varphi\in\Z^{V}$, the switched voltage $\alpha+\partial^{t}\varphi$ has pencil
$\operatorname{diag}(t^{\varphi_u})\,D(t)\,\operatorname{diag}(t^{-\varphi_u})$, so $f$
depends only on the class $[\alpha]\in H^{1}(Y_0;\Z)$.
\end{lemma}

\begin{proof}
$A(t^{-1})_{uv}=\sum_{e\colon u\to v}t^{-\alpha(e)}=\sum_{e\colon v\to u}t^{\alpha(e)}
=A(t)_{vu}$. Switching multiplies the entry $t^{\alpha(e)}$, $e\colon u\to v$, by
$t^{\varphi_v-\varphi_u}$.
\end{proof}

\begin{theorem}[Exceptional zero of order two]\label{IX:thm:zero}
\begin{enumerate}
\item[(i)] $\Theta(0)=0$ and $\Theta'(0)=0$.
\item[(ii)] $a_2=\tfrac12\Theta''(0)=-\,\kappa(Y_0)\,\|h_\alpha\|^{2}$.
\item[(iii)] The zero has order exactly two if and only if $h_\alpha\neq0$, i.e.\ if and
only if $[\alpha]\neq0$ in $H^{1}(Y_0;\Q)$.
\item[(iv)] If $\alpha$ is the indicator of a single edge $e$, then
$\|h_\alpha\|^{2}=1-\Reff(e)$ and $-a_2=\kappa(Y_0)\bigl(1-\Reff(e)\bigr)$ is the number
of spanning trees of $Y_0$ not containing $e$.
\end{enumerate}
\end{theorem}

\begin{proof}
Write $s=T$ and expand $D(1+s)=\Delta_0+sD_1+s^{2}D_2+O(s^{3})$ with
$D_1=D'(1)=-A'(1)=:-W$ and $D_2=\frac12D''(1)=-\frac12A''(1)$. By
Lemma~\ref{IX:lem:sym}, $A'(1)^{t}=-A'(1)$, so $W$ is antisymmetric, and
$(W\mathbf1)_u=\sum_{o(e)=u}\alpha(e)-\sum_{t(e)=u}\alpha(e)=-(\partial\alpha)_u$, i.e.\
$W\mathbf1=-\partial\alpha$. Also
$\mathbf1^{t}A''(1)\mathbf1=\sum_{e\in E^+}\bigl[\alpha(\alpha-1)+\alpha(\alpha+1)\bigr]
=2\|\alpha\|^{2}$, so $\mathbf1^{t}D_2\mathbf1=-\|\alpha\|^{2}$.

Decompose $\Q^{n}=\Q\mathbf1\oplus\mathbf1^{\perp}$ orthogonally, with unit vector
$\mathbf1/\sqrt n$. In this decomposition $\Delta_0=0\oplus\Delta'$ with $\Delta'$
invertible and $\det\Delta'=\prod_{\mu\neq0}\mu=n\kappa_0$ [Ch.~\ref{chap:III}, eq.~(2.2)]. Writing
$E:=sD_1+s^{2}D_2+O(s^{3})$ in block form $\begin{psmallmatrix}e_{00}&e_{0*}\\
e_{*0}&E_{**}\end{psmallmatrix}$, the Schur complement gives
\[
f(1+s)=\det(\Delta'+E_{**})\cdot\bigl(e_{00}-e_{0*}(\Delta'+E_{**})^{-1}e_{*0}\bigr).
\]
Here $e_{00}=\frac1n\mathbf1^{t}E\mathbf1=\frac{s}{n}\mathbf1^{t}D_1\mathbf1
+\frac{s^{2}}n\mathbf1^{t}D_2\mathbf1+O(s^{3})=-\frac{s^{2}}n\|\alpha\|^{2}+O(s^{3})$,
because $\mathbf1^{t}W\mathbf1=0$ by antisymmetry; and $e_{0*},e_{*0}=O(s)$ with
$e_{*0}=\frac{s}{\sqrt n}D_1\mathbf1|_{\mathbf1^\perp}+O(s^{2})$,
$e_{0*}=\frac{s}{\sqrt n}(D_1^{t}\mathbf1)^{t}|_{\mathbf1^\perp}+O(s^{2})$; note
$D_1\mathbf1=-W\mathbf1=\partial\alpha\in\mathbf1^{\perp}$ and
$D_1^{t}\mathbf1=W\mathbf1=-\partial\alpha$. Since $(\Delta'+E_{**})^{-1}=\Delta_0^{+}
+O(s)$ on $\mathbf1^{\perp}$,
\[
e_{0*}(\Delta'+E_{**})^{-1}e_{*0}
=\frac{s^{2}}n(-\partial\alpha)^{t}\Delta_0^{+}(\partial\alpha)+O(s^{3}) .
\]
Hence
$f(1+s)=\bigl(n\kappa_0+O(s)\bigr)\cdot\frac{s^{2}}{n}\bigl(-\|\alpha\|^{2}
+(\partial\alpha)^{t}\Delta_0^{+}(\partial\alpha)\bigr)+O(s^{3})
=-\kappa_0\|\Pi\alpha\|^{2}s^{2}+O(s^{3})$, which is (i) and (ii). (iii): $\Pi\alpha=0$
iff $\alpha$ lies in the cut space $\operatorname{im}\partial^{t}\otimes\Q$ iff
$[\alpha]=0$ in $H^{1}(Y_0;\Q)$. (iv): for $\alpha=e$,
$\|\Pi e\|^{2}=1-e^{t}\partial^{t}\Delta_0^{+}\partial e=1-\Reff(e)$, and
$\kappa_0\Reff(e)$ is the number of spanning trees containing $e$ (Kirchhoff). On $K_4$:
$\Reff=\frac12$, $16\cdot\frac12=8$ trees avoid the chord, and $a_2=-8$ in both
$(1,0,0)$ towers (Table~\ref{IX:tab:battery}).
\end{proof}

\begin{definition}[Graph $\cL$-invariant]\label{IX:def:L}
$\displaystyle\cL(Y_\bullet):=\frac{\Theta''(0)}{2\,\kappa(Y_0)}=-\|h_\alpha\|^{2}
=-\|\Pi_{\ker\mathsf D_p}\alpha\|^{2}\ \in\Q_{\le0}$, nonzero iff $[\alpha]\neq0$
rationally. It depends only on $[\alpha]\in H^{1}(Y_0;\Z)$ and not on $\ell$: the same
rational number serves every prime. On $K_4$ the four towers of [Ch.~\ref{chap:II}, Table~2] have
$\cL=-\frac12,\,-1,\,-\frac32,\,-\frac12$ for $(1,0,0)$, $(1,1,1)$, $(1,2,0)$, $(1,0,0)$
respectively.
\end{definition}

\begin{theorem}[The graph Mazur--Tate--Teitelbaum formula]\label{IX:thm:mtt}
$\displaystyle\tfrac12\Theta''(0)=\cL(Y_\bullet)\cdot\kappa(Y_0)
=\cL(Y_\bullet)\cdot\frac1n\prod_f\#E_f(\F_p)$: the leading coefficient of the graph
$\ell$-adic $L$-function at its exceptional zero is the $\cL$-invariant times the
algebraic $L$-value.
\end{theorem}

\begin{proof}
Theorem~\ref{IX:thm:zero}(ii), Definition~\ref{IX:def:L}, and [Ch.~\ref{chap:I}, Thm.~3.5] for the product
formula for $\kappa(Y_0)$.
\end{proof}

\begin{remark}[The analogy, stated as an analogy]\label{IX:rem:analogy}
For an elliptic curve with split multiplicative reduction at $p$, the $p$-adic
$L$-function has an exceptional zero at $s=1$ and \cite{MTT} conjectured, \cite{GS}
proved, $L_p'(E,1)=\cL(E)\,L(E,1)/\Omega_E$ with $\cL(E)=\log_pq_E/\ord_pq_E$ the
$\cL$-invariant of the Tate period. Theorem~\ref{IX:thm:mtt} has that shape: a zero forced
by a symmetry, an algebraic part that is a product of Frobenius point counts, and a
transcendental-looking factor that is in fact a period-type quantity --- here the Hodge
energy of the class $[\alpha]$, i.e.\ its length in the harmonic metric of $Y_0$. The
zero is of order two rather than one because the graph pencil is self-dual under
$t\mapsto t^{-1}$ (Lemma~\ref{IX:lem:sym}), which kills the first derivative; the analogue
in the classical setting would be a functional equation forcing an even order. We
claim the shape and nothing more.
\end{remark}

\section{The $\ell$-adic bookkeeping}\label{IX:sec:book}

Fix a prime $\ell$ with the tower $Y_k$ connected for all $k$ (equivalently
$\alpha\not\equiv0$ in $H^{1}(Y_0;\F_\ell)$, true in the four towers). Write
$h(T):=g(T)/T^{2}\in\Z[T]$, so $h(0)=a_2$, and let $h=\ell^{\mu}P(T)U(T)$ be its
Weierstrass decomposition in $\La$: $P$ distinguished of degree $\lambda_h$, $U$ a unit
\cite{Wa}. Thus $\mu=\min_i\ord_\ell(\text{coefficients of }h)$ and $\lambda_h$ is the
index of the first coefficient attaining it; $[Ch.~\ref{chap:II}, Rem.~5.4]$ and $[Ch.~\ref{chap:IV}, Cor.~2.3]$ give
the Iwasawa invariants of the tower as $\mu$ and $\lambda=\lambda(g)-1=\lambda_h+1$.
For $j\ge1$ put
\[
N_j:=\ord_\ell\prod_{\zeta\ \mathrm{prim.}\ \ell^{j}}h(\zeta-1)
=\ord_\ell\Res_T\bigl(\Phi_{\ell^{j}}(1+T),\,h(T)\bigr),
\qquad
\varepsilon_j:=N_j-\bigl(\mu\,\varphi(\ell^{j})+\lambda_h\bigr),
\]
the exceptional low-level defects of [Ch.~\ref{chap:III}, \S3].

\begin{lemma}[Character product]\label{IX:lem:char}
For every $k\ge0$,
$\kappa(Y_k)=\kappa(Y_0)\,\ell^{-k}\prod_{\zeta^{\ell^{k}}=1,\ \zeta\neq1}\det D(\zeta)$
{\rm[Ch.~\ref{chap:II}, \S5]}, and consequently
\[
\ord_\ell\kappa(Y_k)=\ord_\ell\kappa_0-k+\sum_{j=1}^{k}\bigl(2+N_j\bigr).
\]
\end{lemma}

\begin{proof}
$\det D(\zeta)=\zeta^{-d}g(\zeta-1)$ and $\prod_{\zeta\neq1}\zeta^{-d}$ is a root of
unity; $g(\zeta-1)=(\zeta-1)^{2}h(\zeta-1)$ and
$\prod_{\zeta\ \mathrm{prim.}\ \ell^{j}}(\zeta-1)=\pm\Phi_{\ell^{j}}(1)=\pm\ell$, of
valuation $1$; group the nontrivial roots by exact order $\ell^{j}$, $1\le j\le k$.
\end{proof}

\begin{theorem}[The bookkeeping identity of {\rm[Ch.~\ref{chap:III}, \S3]}]\label{IX:thm:book}
Let $\rho_1,\dots,\rho_{\lambda_h}$ be the roots of $P$ (all of positive valuation) and
let $j_0\ge1$ be least with $1/\varphi(\ell^{j_0})<\min_i\ord_\ell\rho_i$ ($j_0=1$ if
$\lambda_h=0$). Then:
\begin{enumerate}
\item[(i)] $\varepsilon_j=0$ for all $j\ge j_0$; the defects are finitely supported;
\item[(ii)] for all $k\ge j_0-1$,
\[
\ord_\ell\kappa(Y_k)=\mu\,\ell^{k}+(\lambda_h+1)\,k+\nu_0,\qquad
\nu_0=\ord_\ell\kappa(Y_0)-\mu+\sum_{j\ge1}\varepsilon_j ;
\]
\item[(iii)] in particular $\lambda=\lambda_h+1=\lambda(g)-1$, recovering
{\rm[Ch.~\ref{chap:II}, Rem.~5.4]} and {\rm[Ch.~\ref{chap:IV}, Cor.~2.3]}, and the constant term is the identity
recorded in {\rm[Ch.~\ref{chap:III}, \S3]}.
\end{enumerate}
\end{theorem}

\begin{proof}
(i) For $\zeta$ primitive of order $\ell^{j}$, $\ord_\ell(\zeta-1)=1/\varphi(\ell^{j})$;
if this is smaller than every $\ord_\ell\rho_i$ then
$\ord_\ell(\zeta-1-\rho_i)=1/\varphi(\ell^{j})$ for all $i$, so
$\ord_\ell h(\zeta-1)=\mu+\lambda_h/\varphi(\ell^{j})$, and summing over the
$\varphi(\ell^{j})$ primitive roots gives $N_j=\mu\varphi(\ell^{j})+\lambda_h$. (ii) By
Lemma~\ref{IX:lem:char},
$\ord_\ell\kappa(Y_k)=\ord_\ell\kappa_0-k+2k+\sum_{j\le k}(\mu\varphi(\ell^{j})
+\lambda_h+\varepsilon_j)=\ord_\ell\kappa_0+k+\mu(\ell^{k}-1)+\lambda_hk
+\sum_{j\le k}\varepsilon_j$, and for $k\ge j_0-1$ the last sum is the full sum. (iii)
is read off. Machine confirmation (Table~\ref{IX:tab:battery}): the law holds with the
predicted $(\mu,\lambda,\nu_0)=(3,1,1),(0,3,4),(2,3,3),(0,1,0)$ --- the four rows of
[Ch.~\ref{chap:II}, Table~2] --- at every computed level $k\ge2$, and in fact at every $k\ge0$ except
in the $(1,2,0)$ tower, where $(\varepsilon_1,\varepsilon_2)=(-1,+2)$ and the law
starts at $k=2$, exactly as [Ch.~\ref{chap:III}, \S3] reported.
\end{proof}

\begin{corollary}[The exceptional-zero regime]\label{IX:cor:regime}
The following are equivalent: $\lambda_h=0$; $\ord_\ell a_2=\mu$;
$\ord_\ell\cL(Y_\bullet)=\mu-\ord_\ell\kappa(Y_0)$. In this regime $\varepsilon_j=0$ for
every $j\ge1$, the law of Theorem~\ref{IX:thm:book}(ii) holds for every $k\ge0$, and
\[
\nu_0\;=\;-\,\ord_\ell\cL(Y_\bullet):
\]
the constant term of the Iwasawa law is minus the $\ell$-adic valuation of the
$\cL$-invariant. This is the case in both $(1,0,0)$ towers: $\nu_0=1=-\ord_2(-\tfrac12)$
and $\nu_0=0=-\ord_3(-\tfrac12)$.
\end{corollary}

\begin{proof}
$\mu\le\ord_\ell a_2$ always, with equality iff the constant coefficient attains the
minimum iff $\lambda_h=0$; the third condition is the second rewritten. If $\lambda_h=0$
then $P=1$, $\ord_\ell h(\zeta-1)=\mu$ for every $\zeta$, so $N_j=\mu\varphi(\ell^{j})$
and $\varepsilon_j=0$, $j_0=1$; then $\nu_0=\ord_\ell\kappa_0-\mu
=\ord_\ell\kappa_0-\ord_\ell a_2=-\ord_\ell\cL$.
\end{proof}

\begin{theorem}[The $\cL$-defect identity]\label{IX:thm:defect}
In general,
\[
\nu_0+\ord_\ell\cL(Y_\bullet)\;=\;\bigl(\ord_\ell a_2-\mu\bigr)+\sum_{j\ge1}\varepsilon_j
\;=:\;\delta(Y_\bullet,\ell),
\]
with $\ord_\ell a_2-\mu\ge0$ and equality iff $\lambda_h=0$. On the four towers
$\delta=0,4,2,0$.
\end{theorem}

\begin{proof}
Substitute $\ord_\ell\cL=\ord_\ell a_2-\ord_\ell\kappa_0$ into
Theorem~\ref{IX:thm:book}(ii).
\end{proof}

\begin{remark}[A correction to the specification]\label{IX:rem:align}
The specification asked to prove that $\ord_\ell\cL(Y_\bullet)$ ``aligns exactly'' with
$\nu_0-\ord_\ell\kappa(Y_0)+\mu$. By Theorem~\ref{IX:thm:book} the latter quantity is
$\sum_j\varepsilon_j$, which is $0,0,1,0$ on the four towers, whereas
$\ord_\ell\cL=-1,0,-1,0$: there is no such identity. The statements that hold are
Corollary~\ref{IX:cor:regime} and Theorem~\ref{IX:thm:defect}. The specification's other
claim, that $\ord_\ell a_2=\mu$ ``controls the main term in the $\lambda_h=0$ regime'',
is the definition of that regime; its content is the vanishing of all defects and the
formula $\nu_0=-\ord_\ell\cL$.
\end{remark}

\begin{proposition}[The exceptional zero inside the Iwasawa module]\label{IX:prop:module}
Let $X_\infty\cong\La^{n}/D(1+T)\La^{n}$ be the limit module of {\rm[Ch.~\ref{chap:IV}, Thm.~2.1]}, with
$\car(X_\infty)=(\Theta)$. Then $X_\infty/TX_\infty\cong(\Z\oplus\Jac(Y_0))\otimes\Zl$
has $\Zl$-rank one, while $\ord_T\Theta=2$; at the augmentation prime the localization
is
\[
X_\infty\otimes_{\La}\La_{(T)}\;\cong\;\La_{(T)}\big/T^{2}\La_{(T)} ,
\]
cyclic of length two. The exceptional zero is the statement that the Eisenstein
eigenspace of the Iwasawa module is not semisimple.
\end{proposition}

\begin{proof}
The first isomorphism is the exact control theorem [Ch.~\ref{chap:IV}, Cor.~2.2] at $k=0$.
$\La_{(T)}$ is a discrete valuation ring with uniformizer $T$ and residue field $\Ql$,
so $D(1+T)$ has elementary divisors over it; reducing modulo $T$ gives
$D(1)=\Delta_0$ of rank $n-1$ over $\Ql$, so exactly one elementary divisor is divisible
by $T$, and since the $T$-adic order of the determinant is $2$
(Theorem~\ref{IX:thm:zero}) that divisor is $T^{2}$ times a unit.
\end{proof}

\section{The operator-algebraic formulation}\label{IX:sec:oa}

Setting from [Ch.~\ref{chap:I}]. $\cG_\Gamma$ is the Hecke groupoid, $C^*(\cG_\Gamma)=\cO_{\Anb}$ the
Cuntz--Krieger algebra of the non-backtracking matrix
$\Anb=T^{t}S-J$ [Ch.~\ref{chap:I}, Def.~1.7, Prop.~1.9(ii)], where $S,T\in M_{n\times2m}(\Z)$ are the
start and terminal incidence matrices of oriented edges and $J$ is edge reversal. The
Satake dynamics $\sigma^{(\pi)}$ [Ch.~\ref{chap:I}, Def.~2.1] has partition function
$Z_p(\beta;\pi)=L_p(\beta,\pi_p)$ [Ch.~\ref{chap:I}, Thm.~2.2]; for the Eisenstein parameter
$(\alpha_1,\alpha_2)=(1,q)$ the $\KMS_\beta$ simplex is nonempty exactly at the critical
inverse temperature $\beta=1$, where it is a single Perron state [Ch.~\ref{chap:I}, Thm.~2.4(ii),(iv)];
off the Perron locus the Satake data survive only as $\KMS$ \emph{functionals}
[Ch.~\ref{chap:I}, Thm.~2.5, Rem.~2.6]. Throughout $u=q^{-\beta}$, so $\beta=1$ is $u=q^{-1}$.

\begin{definition}[Character twist]\label{IX:def:twist}
The voltage $\alpha$ is a $\Z$-valued continuous $1$-cocycle $c_\alpha$ on
$\cG_\Gamma$ --- the sum of $\alpha$ along the edges of a Hecke move, well defined on
$(\omega,i-j,\omega')$ because the tails agree --- commuting with the degree cocycle
$c$. It induces a $\mathbb T$-action $\gamma^{\alpha}$ on $C^*(\cG_\Gamma)$ commuting
with $\sigma^{(\pi)}$, and for $t\in\C^{\times}$ the \emph{twisted transfer operator} and
\emph{twisted Hecke operators}
\[
B_t:=B\La_t,\qquad \La_t:=\operatorname{diag}\bigl(t^{\alpha(e)}\bigr)_{e},\qquad
T_k(t)_{uv}:=\sum_{\text{nb paths }u\to v\text{ of length }k}t^{\alpha(\text{path})},
\]
with $B=\Anb$ on $\C^{2m}$; for $|t|=1$ these are the character-twisted Ruelle--Hecke
transfer operator of [Ch.~\ref{chap:I}, Def.~2.1] at the Eisenstein potential and the twisted Hecke
correspondences. Note $T_1(t)=A(t)$ and $\La_tJ\La_t=J$, since $\alpha(\bar e)=-\alpha(e)$.
\end{definition}

\begin{lemma}[Twisted Hecke recursion]\label{IX:lem:rec}
$A(t)T_k(t)=T_{k+1}(t)+qT_{k-1}(t)$ for $k\ge2$, $A(t)^{2}=T_2(t)+(q{+}1)I$, hence
\[
\sum_{k\ge0}T_k(t)u^{k}=(1-u^{2})\bigl(I-uA(t)+qu^{2}I\bigr)^{-1}
\]
as an identity of matrix-valued rational functions of $(u,t)$.
\end{lemma}

\begin{proof}
$A(t)T_k(t)$ sums over an edge $e\colon u\to w$ followed by a non-backtracking path
$w\to v$ of length $k$, with weight $t^{\alpha(e)}t^{\alpha(\text{path})}$. The terms in
which the path begins with $\bar e$ have weight $t^{\alpha(e)+\alpha(\bar e)}
t^{\alpha(\text{rest})}=t^{\alpha(\text{rest})}$, the rest being a non-backtracking path
$u\to v$ of length $k-1$ whose first edge is not $e$; for each such rest there are $q$
admissible $e$. This is the proof of [Ch.~\ref{chap:I}, Lem.~1.5] with the multiplicative weight
carried along; the generating function follows by the same algebra.
\end{proof}

\begin{theorem}[Twisted Bass identity]\label{IX:thm:bass}
With $P(u,t):=\det\bigl(I_n-uA(t)+qu^{2}I_n\bigr)$,
\[
\det\bigl(I_{2m}-uB\La_t\bigr)=(1-u^{2})^{\,r-1}\,P(u,t) .
\]
\end{theorem}

\begin{proof}
Recall $S\La_tT^{t}=A(t)$ (entry $(u,v)$ sums $t^{\alpha(e)}$ over $e\colon u\to v$),
$SJ=T$, $TT^{t}=(q{+}1)I$, and $\La_tJ\La_t=J$. The operator $J\La_t$ is an involution
($(J\La_t)^{2}=J\La_tJ\La_t=JJ=I$) with trace $0$ (no edge is its own reverse), so its
eigenvalues are $\pm1$ with multiplicity $m$ each and
$\det(I+uJ\La_t)=(1-u^{2})^{m}$, with $(I+uJ\La_t)^{-1}=(1-u^{2})^{-1}(I-uJ\La_t)$.
Then $I-uB\La_t=(I+uJ\La_t)-uT^{t}S\La_t$, and by $\det(I-XY)=\det(I-YX)$,
\begin{align*}
\det(I-uB\La_t)
&=(1-u^{2})^{m}\det\Bigl(I_n-\tfrac{u}{1-u^{2}}\,S\La_t(I-uJ\La_t)T^{t}\Bigr)\\
&=(1-u^{2})^{m-n}\det\Bigl((1-u^{2})I_n-u\bigl(A(t)-uSJT^{t}\bigr)\Bigr)\\
&=(1-u^{2})^{m-n}\det\bigl(I-uA(t)+qu^{2}I\bigr),
\end{align*}
using $S\La_tJ\La_tT^{t}=SJT^{t}=TT^{t}=(q{+}1)I$ and $m-n=r-1$. For $t=1$ this is
Bass's identity [Ch.~\ref{chap:I}, Prop.~3.4]; for $|t|=1$ it is the Artin--Ihara $L$-function formula
of \cite{ST}. Verified as a polynomial identity in $(u,t)$ on all four voltage
assignments (Table~\ref{IX:tab:battery}).
\end{proof}

\begin{theorem}[Exceptional zero at the critical temperature]\label{IX:thm:crit}
Let $P(u,t)$ be as above and $Z(\beta;t):=\tr\sum_kT_k(t)q^{-k\beta}$ the twisted
partition function of the Eisenstein Satake dynamics.
\begin{enumerate}
\item[(i)] Functional equation: $P\bigl(\tfrac1{qu},t\bigr)=(qu^{2})^{-n}P(u,t)$.
\item[(ii)] At the critical inverse temperature $\beta=1$,
$P(q^{-1},t)=q^{-n}\,\Theta(t-1)$: the twisted Bass determinant of the Satake dynamics
at $\beta=1$ is the graph $\ell$-adic $L$-function.
\item[(iii)] At $(u,t)=(q^{-1},1)$ the zero is simple in the temperature direction and
of order two in the character direction:
\begin{gather*}
\partial_uP(q^{-1},1)=-\,n(q-1)\,q^{1-n}\,\kappa(Y_0),\\
\partial_tP(q^{-1},1)=0,\qquad
\tfrac12\partial_t^{2}P(q^{-1},1)=q^{-n}\,\cL(Y_\bullet)\,\kappa(Y_0) .
\end{gather*}
The ratio $\tfrac12\partial_t^{2}P/\partial_uP=-\cL(Y_\bullet)/(nq(q-1))
=\|h_\alpha\|^{2}/(nq(q-1))$ is independent of $\kappa(Y_0)$.
\item[(iv)] The unique $\KMS_1$ state of {\rm[Ch.~\ref{chap:I}, Thm.~2.4]} is the $t=1$ member of the
family; for $t\neq1$ on the unit circle with $t^{\alpha(C)}\neq1$ for some cycle $C$
there is no $\KMS_1$ state, only the functional of {\rm[Ch.~\ref{chap:I}, Rem.~2.6]}, and the critical
partition function has a double pole at the trivial character with
\[
Z(1;t)=q\,(1-q^{-2})\,\tr D(t)^{-1}
=\frac{n\,(q-q^{-1})}{\cL(Y_\bullet)}\,\frac1{(t-1)^{2}}+O\bigl((t-1)^{-1}\bigr).
\]
\end{enumerate}
\end{theorem}

\begin{proof}
(i) $I-\frac1{qu}A+\frac{q}{q^{2}u^{2}}I=(qu^{2})^{-1}(qu^{2}I-uA+I)$. (ii) At $u=q^{-1}$,
$I-q^{-1}A(t)+q^{-1}I=q^{-1}D(t)$. (iii) $\partial_tP(q^{-1},t)=q^{-n}f'(t)$ and
$\frac12\partial_t^{2}P(q^{-1},1)=q^{-n}a_2$, so Theorem~\ref{IX:thm:zero} gives the
$t$-derivatives. For the $u$-derivative, $\partial_u\det M(u)=\tr\bigl(\adj M(u)\,M'(u)\bigr)$
with $M(q^{-1})=q^{-1}\Delta_0$, $\adj(q^{-1}\Delta_0)=q^{1-n}\adj\Delta_0
=q^{1-n}\kappa_0\mathbf1\mathbf1^{t}$ (all cofactors of a Laplacian equal $\kappa_0$),
and $M'(q^{-1})=-T_p+2I$, so $\partial_uP=q^{1-n}\kappa_0\mathbf1^{t}(2I-T_p)\mathbf1
=q^{1-n}\kappa_0n(2-(q{+}1))$. (iv) Existence of a $\KMS_1$ state requires
$\rho(q^{-1}B\La_t)=1$ [Ch.~\ref{chap:I}, Thm.~2.4(i)]; for $|t|=1$ the twisted Laplacian $D(t)$ is
Hermitian positive semidefinite with $D(t)v=0$ forcing $|v|$ constant and
$t^{\alpha(C)}=1$ on every cycle, so under the hypothesis $D(t)$ is positive definite,
$P(q^{-1},t)\neq0$, and the pressure equation fails (Perron obstruction, as in
[Ch.~\ref{chap:I}, Thm.~2.5]). By Lemma~\ref{IX:lem:rec},
$Z(1;t)=(1-q^{-2})\tr\bigl(I-q^{-1}A(t)+q^{-1}I\bigr)^{-1}=q(1-q^{-2})\tr D(t)^{-1}$.
The Perron eigenvalue branch of $D(t)$ through $0$ at $t=1$ is
$\lambda(1+s)=-\frac{\|h_\alpha\|^{2}}{n}s^{2}+O(s^{3})$ (this is the Schur-complement
computation in the proof of Theorem~\ref{IX:thm:zero}: $f(1+s)/\det\Delta'$), and the other
$n-1$ eigenvalues are bounded away from zero, so $\tr D(t)^{-1}=
-\frac{n}{\|h_\alpha\|^{2}}(t-1)^{-2}+O((t-1)^{-1})=\frac{n}{\cL(Y_\bullet)}(t-1)^{-2}
+O((t-1)^{-1})$. On $K_4$ the leading coefficients are $-12,-6,-4,-12$
(Table~\ref{IX:tab:battery}).
\end{proof}

\begin{remark}[Which derivative]\label{IX:rem:whichderivative}
The specification asked for ``the second derivative of the KMS functional at $\beta=1$''.
Theorem~\ref{IX:thm:crit} locates the derivatives precisely: at the critical point the
\emph{temperature} derivative is simple and carries the algebraic $L$-value
$\kappa(Y_0)$; the \emph{character} derivative is of second order and carries
$\cL(Y_\bullet)\kappa(Y_0)$; their ratio is the $\cL$-invariant up to the explicit
normalization $nq(q-1)$, with $\kappa(Y_0)$ cancelling --- which is the Greenberg--Stevens
shape \cite{GS}, a ratio of a derivative in the ``weight'' direction to one in the
cyclotomic direction. The specification's request for a cyclic-cohomology formulation
is not met here: the natural candidate pairs the derivation generating $\gamma^{\alpha}$
with the Fredholm module of [Ch.~\ref{chap:II}, Thm.~3.3], and we have not established that pairing.
We record it as open (Remark~\ref{IX:rem:scope}) rather than dress the second derivative in
cohomological language it has not earned.
\end{remark}

\begin{theorem}[The abelian shadow limit carries $\Theta$]\label{IX:thm:shadow}
Let $(Y_k)$ be a connected abelian $\Z/\ell^{k}$ voltage tower, $\cO_k:=\cO_{B(Y_k)}$ the
shadow algebras of {\rm[Ch.~\ref{chap:VI}]} ($c=q$ at every level, all levels being $(q{+}1)$-regular),
and $\pi\colon Y_{k+1}\to Y_k$ the covering maps. Then $\pi$ extends to an in-covering
$\pi^{\sharp}$ of shadow graphs, {\rm[Ch.~\ref{chap:VIII}, Lem.~2.1]} yields unital injective
$*$-homomorphisms $\varphi_k\colon\cO_k\to\cO_{k+1}$, and
$\cO_\infty^{\mathrm{ab}}:=\varinjlim(\cO_k,\varphi_k)$ is a unital Kirchberg algebra with
\[
K_1(\cO_\infty^{\mathrm{ab}})\cong\Z,\qquad
\Tor K_0(\cO_\infty^{\mathrm{ab}})\cong\Hom_{\mathrm{cont}}\Bigl(\varprojlim_k\bigl(\Jac(Y_k),\pi_*\bigr),\Q/\Z\Bigr),
\qquad K_0/\Tor\cong\Z[1/\ell].
\]
The deck transformations induce a continuous action of $\Zl=\varprojlim\Z/\ell^{k}$ on
$\cO_\infty^{\mathrm{ab}}$ commuting with the gauge action, and the $\ell$-primary torsion
of $K_0$ is, as a discrete $\La$-module, the Pontryagin dual of the compact $\La$-module
$\Tor X_\infty$ of {\rm[Ch.~\ref{chap:IV}, Thm.~2.1]}, whose characteristic ideal is $(\Theta)$. The
ideal $(\Theta)$ is stable under the involution $1+T\mapsto(1+T)^{-1}$
(Lemma~\ref{IX:lem:sym}), so the same $\Theta$ --- with its exceptional zero and its
$\cL$-invariant --- is the characteristic element on both sides of the duality.
\end{theorem}

\begin{proof}
A covering map is bijective on in-arcs and on out-arcs at every vertex, so
$\pi^{\sharp}$ ($v\mapsto\pi(v)$, $v'\mapsto\pi(v)'$, shadow arcs label-wise) is an
in-covering, and [Ch.~\ref{chap:VIII}, Lem.~2.1, Thm.~2.3] apply verbatim with $\tau$ replaced by
$\pi$: $K_0(\varphi_k)$ is the pullback $\pi^{*}$ on $\coker\Delta_k$ (the Laplacians
are symmetric, so the transpose convention of [Ch.~\ref{chap:VIII}, Conv.~1.1] is invisible), which is
the Pontryagin adjoint of the norm $\pi_*$ of [Ch.~\ref{chap:IV}, Lem.~1.2--1.3] by
[Ch.~\ref{chap:VIII}, Lem.~1.5]; $\pi^{*}\mathbf1=\mathbf1$ gives unitality, $\deg\pi^{*}=\ell\deg$
gives the free part, and $K_1(\varphi_k)=\mathrm{id}$ as in [Ch.~\ref{chap:VIII}, Thm.~2.3(iii)]. The
limit is Kirchberg by [Ch.~\ref{chap:VIII}, Thm.~3.2] and its $K$-theory follows as in
[Ch.~\ref{chap:VIII}, Thm.~3.3]. Each deck generator $\sigma_k$ is an automorphism of $Y_k^{\sharp}$
covering $\sigma_{k-1}$, hence induces an automorphism of $\cO_k$ with
$\sigma_{k+1}\circ\varphi_k=\varphi_k\circ\sigma_k$ (both send $p_v$ to
$\sum_{\pi(w)=\sigma_kv}p_w$), so $\Zl$ acts on the limit, continuously because on each
$\cO_k$ the action factors through $\Z/\ell^{k}$. On $K_0(\cO_k)=\coker\Delta_k\otimes\Zl$
the deck action is multiplication by $\sigma\in\Zl[G_k]$ [Ch.~\ref{chap:IV}, Lem.~1.1], so the
$\ell$-part of $\Tor K_0(\cO_\infty^{\mathrm{ab}})=\varinjlim\Jac(Y_k)^{\vee}$ is the
Pontryagin dual of $\varprojlim\Jac(Y_k)\otimes\Zl=\Tor X_\infty$ with the
contragredient $\La$-action, and $\car(X_\infty)=(\Theta)$ is [Ch.~\ref{chap:IV}, Thm.~2.1]. The
involution stability is $f(t^{-1})=f(t)$.
\end{proof}

\begin{remark}[Two algebras, one zero]\label{IX:rem:twoalgebras}
The shadow algebras are Hecke-functorial and not Hecke-dynamical [Ch.~\ref{chap:VI}, Rem.~4.3]: their
gauge dynamics reach the Laplacian at $u=1$, i.e.\ $\beta=0$, where pure infiniteness
forbids traces and hence $\KMS_0$ states. So the exceptional zero has two
operator-algebraic homes, and they do not overlap: the \emph{dynamical} side ---
partition functions, the critical temperature $\beta=1$, the double pole of
Theorem~\ref{IX:thm:crit}(iv) --- lives on $C^*(\cG_\Gamma)$ of [Ch.~\ref{chap:I}]; the
\emph{$K$-theoretic} side --- $\Theta$ as characteristic element,
Proposition~\ref{IX:prop:module} --- lives on the equivariant limit
$\cO_\infty^{\mathrm{ab}}$ of Theorem~\ref{IX:thm:shadow}. The specification's request to
read the derivative off ``the gauge action of $\cO_\infty$'' conflates the two; by [Ch.~\ref{chap:II}, Thm.~1.6] no single algebra can carry both.
\end{remark}

\section{The Iwahori tower: a simple zero and no $\cL$-invariant}\label{IX:sec:iwahori}

\begin{theorem}[Residue constancy along the Iwahori tower]\label{IX:thm:iwahori}
For the Iwahori path tower $(Y_k)_{k\ge1}$ of {\rm[Ch.~\ref{chap:II}, \S5]} over a $(q{+}1)$-regular
base with $q=p$, let $\zeta_k(u)^{-1}:=\det(I-uA_k)$. Then:
\begin{enumerate}
\item[(i)] $\zeta_k(u)^{-1}$ is independent of $k\ge1$ and equals
$\det(I-u\Anb)=(1-u^{2})^{r-1}P(u,1)$;
\item[(ii)] $u=q^{-1}$ is a simple zero of $\zeta_k(u)^{-1}$ at every level, with
\[
\frac{d}{du}\det(I-uA_k)\Big|_{u=q^{-1}}=-\,q^{\,2-n_k}\,n_k\,\kappa(Y_k);
\]
\item[(iii)] consequently $n_k\,\kappa(Y_k)\,q^{-n_k}$ is independent of $k\ge1$, which is
the recursion $\kappa(Y_{k+1})=q^{\,n_k(q-1)-1}\kappa(Y_k)$ of {\rm[Ch.~\ref{chap:III}, Thm.~2.1]}, and
comparing with level $0$,
\[
\kappa(Y_1)=\frac{(1-q^{-2})^{\,r-1}(q-1)\,n\,q^{\,n_1-1-n}}{n_1}\,\kappa(Y_0);
\]
\item[(iv)] the pseudo-Iwasawa law {\rm[Ch.~\ref{chap:III}, Thm.~2.2]} reads
$\ord_p\kappa(Y_k)=n_k-\ord_p n_k+\mathrm{const}=\mu p^{k}-k+\nu_0$: the Euler term $-k$
is the $p$-adic valuation of the vertex count $n_k=n(q{+}1)q^{k-1}$, i.e.\ of the
Eisenstein normalization in the residue.
\end{enumerate}
\end{theorem}

\begin{proof}
(i) $A_1=\Anb$, and $\det(zI-A_{k+1})=z^{n_{k+1}-n_k}\det(zI-A_k)$ [Ch.~\ref{chap:III}, eq.~(2.1)] gives
$\det(I-uA_{k+1})=\det(I-uA_k)$; the last equality is Bass [Ch.~\ref{chap:I}, Prop.~3.4]. (ii) With
$M(u)=I-uA_k$, $M(q^{-1})=q^{-1}\Delta_k$ has the simple eigenvalue $0$ (Perron), so the
zero is simple; $\adj(q^{-1}\Delta_k)=q^{1-n_k}\kappa(Y_k)\mathbf1\mathbf1^{t}$ since all
cofactors of the Eulerian Laplacian equal $\kappa(Y_k)$ [Ch.~\ref{chap:III}, proof of Thm.~2.1, Step~2],
and $\tr(\mathbf1\mathbf1^{t}(-A_k))=-\mathbf1^{t}A_k\mathbf1=-qn_k$. (iii) Equate the
derivatives at consecutive levels, using $n_{k+1}=qn_k$; for the level-$0$ formula
differentiate $(1-u^{2})^{r-1}P(u,1)$ at $q^{-1}$ using Theorem~\ref{IX:thm:crit}(iii). (iv)
$\ord_p(n_k\kappa_kq^{-n_k})$ constant gives $\ord_p\kappa_k=n_k-\ord_pn_k+\mathrm{const}$,
and $\ord_pn_k=\ord_p(n(q{+}1))+k-1$. On $K_4$: $n_k\kappa_kq^{-n_k}=\frac98$ at
$k=1,2,3$, the derivative is $-\frac92$ at every level, and
$\kappa(Y_1)=\frac{(3/4)^{2}\cdot4\cdot2^{7}}{12}\cdot16=384$ (Table~\ref{IX:tab:battery}).
\end{proof}

\begin{remark}[No $\cL$-invariant for the non-normal tower]\label{IX:rem:noL}
The specification asked for $\Theta''(0)$, the $\cL$-invariant and the defects
$\varepsilon_j$ ``on the Iwahori path tower as well''. There is no $\Theta$ there: the
tower carries no deck action, hence no character variable, hence no pencil, and the
Hecke variable $u$ is the only direction available. In that direction the zero at the
critical temperature is simple at every level (Theorem~\ref{IX:thm:iwahori}(ii)), with
residue $\kappa(Y_k)$ rather than a second-order coefficient. What replaces the
$\cL$-invariant is the tower constant $n_k\kappa(Y_k)q^{-n_k}$, and what replaces the
bookkeeping identity is the observation that its constancy \emph{is} the growth law.
This gives the Euler term $-k$ its third reading, after the Eisenstein rank removed
once per level [Ch.~\ref{chap:III}, Thm.~2.4] and the Eisenstein index through the connecting map of
the snake [Ch.~\ref{chap:VII}, Cor.~2.2]: it is the $p$-adic valuation of the vertex count, the
normalization $\prod_{\mu\neq0}\mu=n_k\kappa_k$ of [Ch.~\ref{chap:III}, eq.~(2.2)] read $p$-adically.
The three readings are one fact seen from three sides, and none of them is an
$\cL$-invariant.
\end{remark}

\section{Scope, residue, corrections}\label{IX:sec:scope}

\begin{remark}[Solved, open, impossible]\label{IX:rem:scope}
\emph{Solved:} the direction of [Ch.~\ref{chap:III}, \S3], made precise. The double zero is a theorem
with the closed form $\frac12\Theta''(0)=-\kappa(Y_0)\|h_\alpha\|^{2}$; the $\cL$-invariant
is defined, computed, and shown to be a Hodge energy and $\ell$-independent; the
bookkeeping identity holds with an exact description of the defects, and the constant
term is $-\ord_\ell\cL$ in the exceptional-zero regime; the twisted Bass identity and
the critical-temperature theorem give the Mazur--Tate--Teitelbaum and
Greenberg--Stevens shapes on $C^*(\cG_\Gamma)$; the abelian shadow limit carries
$\Theta$ as equivariant characteristic element. \emph{Open:} (a) a cyclic-cohomology
formulation of Theorem~\ref{IX:thm:crit}(iii) (Remark~\ref{IX:rem:whichderivative}); (b) a
converse to Theorem~\ref{IX:thm:defect} --- whether $\delta(Y_\bullet,\ell)=0$ forces
$\lambda_h=0$; (c) a structural formula for the defects $\varepsilon_j$ at $j<j_0$ in
terms of the Newton polygon of $P$ at the levels $\ell^{j}$, beyond the resultant
definition; (d) everything about the non-normal tower that [Ch.~\ref{chap:V}, Prob.~5.1] and
[Ch.~\ref{chap:VII}, Rem.~3.1] left open, which this paper does not touch. \emph{Impossible or
false as stated:} an $\cL$-invariant for the Iwahori tower in the sense of
Definition~\ref{IX:def:L} (Remark~\ref{IX:rem:noL}); the alignment formula of the
specification (Remark~\ref{IX:rem:align}); a single algebra carrying both the dynamical
and the $K$-theoretic side (Remark~\ref{IX:rem:twoalgebras}, by [Ch.~\ref{chap:II}, Thm.~1.6]).
\end{remark}

\begin{remark}[Five corrections to the task specification]\label{IX:rem:corrections}
(1) $\Theta(T)=\det D(1+T)$ lies in $\Z[T,(1+T)^{-1}]\subset\La$, not in $\Zl[T]$; the
cleared $g$ is the polynomial, and $a_2$ is the same for both. (2)
``$\ord_\ell a_2=\mu$ in the $\lambda_h=0$ regime'' is the definition of the regime;
the theorem is Corollary~\ref{IX:cor:regime}. (3) The alignment
``$\ord_\ell\cL\leftrightarrow\nu_0-\ord_\ell\kappa_0+\mu$'' is false;
Theorem~\ref{IX:thm:defect} is the identity that holds. (4) The KMS second derivative is in
the character direction, at $\beta=1$, on $C^*(\cG_\Gamma)$ and not on the shadow limit;
the temperature derivative carries $\kappa(Y_0)$ (Theorem~\ref{IX:thm:crit},
Remarks~\ref{IX:rem:whichderivative}--\ref{IX:rem:twoalgebras}); no cyclic-cohomology
statement is claimed. (5) The Iwahori tower has no $\Theta$, no $\Theta''$, and no
$\varepsilon_j$; Theorem~\ref{IX:thm:iwahori} is what it has. Separately, [Ch.~\ref{chap:III}, Rem.~1.11]
and [Ch.~\ref{chap:VIII}] cite the tail-groupoid $K$-theory as [Ch.~\ref{chap:I}, Prop.~1.7(iii)]; in the compiled
numbering of [Ch.~\ref{chap:I}] it is Proposition~1.9(iii), and [Ch.~\ref{chap:VIII}] has been corrected accordingly.
\end{remark}

\section{The verification battery}\label{IX:sec:battery}

Table~\ref{IX:tab:battery} lists the exact checks of \texttt{exceptional\_zero.py}; the
letters match the docstring of the script and the references in the text. The four
abelian towers are those of [Ch.~\ref{chap:II}, Table~2]; the growth laws were recomputed directly
from the reduced Laplacians of the derived graphs ($k\le5$ for $\ell=2$, up to $128$
vertices; $k\le3$ for $\ell=3$) and independently from the character product of
Lemma~\ref{IX:lem:char}, and the two agree at every level.

\begin{table}[ht]
\centering\scriptsize
\renewcommand{\arraystretch}{1.2}
\begin{tabular}{lcl}
\toprule
check & towers / levels & result\\
\midrule
(Z) $\Theta(0)=\Theta'(0)=0$; $a_2=-8,-16,-24,-8$; $g=-8T^{2}$, $-T^{4}-16T^{3}-16T^{2}$, $-4T^{4}-24T^{3}-24T^{2}$ & 4 towers & exact\\
(H) $a_2=-\kappa_0\|h_\alpha\|^{2}$; $\|h_\alpha\|^{2}=\frac12,1,\frac32,\frac12$; $\cL=-\frac12,-1,-\frac32,-\frac12$ & 4 towers & exact\\
(H$'$) single chord: $\Reff=\frac12$, $8$ spanning trees avoid the chord $=-a_2$ & $(1,0,0)$ & exact\\
(K) double pole of $Z(1;t)$: leading coefficient $=n(q-q^{-1})/\cL=-12,-6,-4,-12$ & 4 towers & exact\\
(M) Newton data $(\mu,\lambda_h)=(3,0),(0,2),(2,2),(0,0)$; regime iff $\ord_\ell a_2=\mu$ & 4 towers & exact\\
(E) $\varepsilon_j$, $j\le5$: all zero except $(\varepsilon_1,\varepsilon_2)=(-1,2)$ for $(1,2,0)$ & 4 towers & exact\\
(G) $\ord_\ell\kappa(Y_k)$ direct $=$ character formula $=$ law; $(\mu,\lambda,\nu_0)$ as [Ch.~\ref{chap:II}, Table~2] & $k\le5$ / $k\le3$ & exact\\
(L) $\nu_0+\ord_\ell\cL=(\ord_\ell a_2-\mu)+\sum\varepsilon_j=0,4,2,0$; regime $\Rightarrow\nu_0=-\ord_\ell\cL$ & 4 towers & exact\\
(B) $\det(I-uB\La_t)=(1-u^{2})^{r-1}P(u,t)$ as a polynomial identity in $(u,t)$ & 4 towers & exact\\
(F) functional equation; $P(q^{-1},t)=q^{-n}\Theta(t-1)$; $\partial_uP=-8$; $\frac12\partial_t^{2}P=q^{-n}a_2$ & 4 towers & exact\\
(S) $\rk D(1)=3=n-1$, $\ord_T\Theta=2$ & 4 towers & exact\\
(I) Iwahori: level-independent $\det(I-uA_k)$; $n_k\kappa_kq^{-n_k}=\frac98$; slope $-\frac92$ at $u=\frac12$; $\kappa(Y_1)=384$ & $k\le3$ & exact\\
\bottomrule
\end{tabular}
\medskip
\caption{Exact verification (\texttt{exceptional\_zero.py}) on $K_4$ ($q=2$, $n=4$,
$r=3$, $\kappa_0=16$). Tower order: $(1,0,0),\ell=2$; $(1,1,1),\ell=2$;
$(1,2,0),\ell=2$; $(1,0,0),\ell=3$. For the fixed towers each line is itself a proof;
the general statements are proved in the text.}
\label{IX:tab:battery}
\end{table}

\section*{Acknowledgement of provenance and caveats}

Unrefereed preprint. Every numbered statement is proved from Papers I--VIII, finite
linear algebra, Weierstrass preparation \cite{Wa} and the character product formula;
the twisted Bass identity is proved rather than quoted from \cite{ST}. The
Mazur--Tate--Teitelbaum and Greenberg--Stevens references \cite{MTT,GS} fix the shape of
an analogy (Remarks~\ref{IX:rem:analogy}, \ref{IX:rem:whichderivative}) and are load-bearing
for nothing. Remark~\ref{IX:rem:corrections} records the corrections to the specification
that prompted the paper and one citation correction propagated from [Ch.~\ref{chap:III}] to [Ch.~\ref{chap:VIII}].
Bibliographic data of \cite{MTT,GS,ST,Ba} were checked against publisher or indexing
pages (titles, venues, years, volumes, pages); \cite{Wa} is the entry verified in [Ch.~\ref{chap:IV}].
